% !TeX program = xelatex
\ifdefined\BookMaster\else
\documentclass[11pt,oneside]{article}
% !TeX program = xelatex
\usepackage[b5paper,left=22mm,right=22mm,top=20mm,bottom=22mm,footskip=11mm]{geometry}
\usepackage{fontspec}
\setmainfont{TeX Gyre Termes}
\setsansfont{TeX Gyre Heros}
\usepackage{amsmath,amssymb,mathtools,amsthm}
\renewcommand{\cdots}{\mathinner{\cdotp\cdotp\cdotp}}
\usepackage{graphicx,adjustbox,array,booktabs,longtable,multirow,enumitem,changepage,needspace}
\usepackage[table]{xcolor}
\usepackage{fancyhdr}
\raggedbottom
\interfootnotelinepenalty=10000
\newcommand{\BookFootnote}[2]{\begingroup\renewcommand{\thefootnote}{\fnsymbol{footnote}}\footnote[#1]{#2}\endgroup}
\newcommand{\BookFootnoteMark}[1]{\begingroup\renewcommand{\thefootnote}{\fnsymbol{footnote}}\footnotemark[#1]\endgroup}
\newcommand{\BookFootnoteText}[2]{\begingroup\renewcommand{\thefootnote}{\fnsymbol{footnote}}\footnotetext[#1]{#2}\endgroup}
\newenvironment{BookFloat}[1]{\par\addvspace{.6\baselineskip}\noindent\def\BookFloatID{#1}\begin{minipage}{\linewidth}\phantomsection\label{bookfloat:#1:start}}{\phantomsection\label{bookfloat:\BookFloatID:end}\end{minipage}\par\addvspace{.6\baselineskip}}

\usepackage[unicode,hidelinks]{hyperref}
\pagestyle{fancy}
\fancyhf{}
\newcommand{\BookFooter}{}
\fancyfoot[L]{\parbox[b]{\dimexpr\textwidth-15mm}{\footnotesize\BookFooter}}
\fancyfoot[R]{\footnotesize\thepage}
\renewcommand{\headrulewidth}{0pt}
\renewcommand{\footrulewidth}{0pt}
\setlength{\parindent}{0pt}
\setlength{\parskip}{0.2em}
\setlength{\emergencystretch}{3em}
\setlist{nosep,leftmargin=1.7em,labelsep=.65em}
\newif\ifBookNarrative
\BookNarrativetrue
\newenvironment{BookBody}{\setlength{\parindent}{0pt}}{}
\newcommand{\BookStatement}[1]{\par\addvspace{.6\baselineskip}\noindent #1\enspace\ignorespaces}
\newcommand{\BookLabel}[1]{\par\addvspace{.6\baselineskip}\needspace{3\baselineskip}\noindent\textbf{#1}\quad\ignorespaces}
\newcommand{\BookHeading}[1]{\par\addvspace{\baselineskip}\needspace{4\baselineskip}\noindent{\large\bfseries #1}\par\medskip}
\newcommand{\BookSubheading}[1]{\par\addvspace{.7\baselineskip}\needspace{3\baselineskip}\noindent{\bfseries #1}\par\smallskip}
\newcommand{\BookTextHeading}[1]{\par\addvspace{.7\baselineskip}\needspace{3\baselineskip}\noindent #1\par\smallskip}
\newcommand{\BookSectionTitle}[2]{\par\addvspace{1.2\baselineskip}\needspace{5\baselineskip}\noindent\colorbox{gray!25}{\sffamily\bfseries\small SECTION #1}\quad{\large\bfseries #2}\par\medskip}
\newcommand{\BookExerciseTitle}[1]{\par\addvspace{\baselineskip}\needspace{4\baselineskip}\noindent{\color{gray}\rule{1.5mm}{1.5mm}}\enspace{\large\bfseries Exercises #1}\par\medskip}
\newcommand{\BookBullet}[1]{\begin{list}{\textbullet}{\setlength{\leftmargin}{1.7em}\setlength{\labelwidth}{1em}\setlength{\labelsep}{.7em}\setlength{\topsep}{0pt}\setlength{\parsep}{0pt}\setlength{\itemsep}{0pt}}\item #1\end{list}}
\newcommand{\BookExercise}[2]{\begin{list}{\textbf{#1}}{\setlength{\leftmargin}{2.2em}\setlength{\labelwidth}{1.5em}\setlength{\labelsep}{.7em}\setlength{\topsep}{.15em}\setlength{\parsep}{0pt}\setlength{\itemsep}{0pt}}\item #2\end{list}}
\newcommand{\BookEnd}[1]{\unskip\nobreak\hfill\ensuremath{#1}}
\newcommand{\BookSource}[2]{\phantomsection\label{sourcepage:#1}} % Source-page anchors, retained through reflow.
\usepackage{multicol}
\usepackage{tikz}
\newcommand{\BookPartEntry}[2]{\par\noindent\makebox[17mm][l]{\bfseries Section #1}#2\par}
\newcommand{\BookPartTitle}[3]{\par\bigskip\noindent\begin{tikzpicture}\node[anchor=north west,inner sep=0pt,outer sep=0pt] (bookpart) at (0,0) {\begin{minipage}{\dimexpr\linewidth-.5pt}\hspace*{26mm}\begin{minipage}[t]{\dimexpr\linewidth-26mm}\begin{adjustbox}{max width=\linewidth}{\fontsize{20}{24}\selectfont\bfseries #2}\end{adjustbox}\par\vspace{6mm}\hrule\vspace{6mm}{\sffamily\footnotesize\setlength{\parskip}{.3em}#3}\vspace{3mm}\end{minipage}\end{minipage}};\path[shade,top color=gray!60,bottom color=white] (0,0) rectangle ([xshift=24mm]bookpart.south west);\node[anchor=north,font=\sffamily\bfseries\small] at (12mm,-2mm) {PART};\node[anchor=north,font=\sffamily\bfseries\Huge] at (12mm,-7mm) {#1};\end{tikzpicture}\par\bigskip}
\newsavebox{\BookMathBox}
\makeatletter
\newcommand{\BookInline}[1]{\leavevmode\hskip0pt\relax\begingroup\in@{\begin}{#1}\ifin@\sbox{\BookMathBox}{\(#1\)}\ifdim\wd\BookMathBox>\dimexpr\linewidth-1em\relax\resizebox{\dimexpr\linewidth-1em\relax}{!}{\usebox{\BookMathBox}}\else\(#1\)\fi\else\in@{\left}{#1}\ifin@\sbox{\BookMathBox}{\(#1\)}\ifdim\wd\BookMathBox>\dimexpr\linewidth-1em\relax\resizebox{\dimexpr\linewidth-1em\relax}{!}{\usebox{\BookMathBox}}\else\(#1\)\fi\else\(#1\)\fi\fi\endgroup\hskip0pt\relax}
\makeatother
\DeclareMathOperator{\Aut}{Aut}
\DeclareMathOperator{\End}{End}
\DeclareMathOperator{\Hom}{Hom}
\DeclareMathOperator{\GF}{GF}
\DeclareMathOperator{\irr}{irr}
\DeclareMathOperator{\lt}{lt}
\DeclareMathOperator{\lp}{lp}
\providecommand{\textcircled}[1]{\textnormal{(#1)}}

\begin{document}
\fi
\renewcommand{\BookFooter}{Section 31 Prime and Maximal Ideals}
\pdfbookmark[0]{Section 31 Prime and Maximal Ideals}{31_prime_and_maximal_ideals}
\ifdefined\BookMaster\else\setcounter{page}{250}\fi
% Source PDF page 267; original printed page 250
\BookSource{267}{250}
\BookNarrativetrue

\begin{BookBody}

\BookSectionTitle{31}{Prime and Maximal Ideals}

\noindent Exercises 12 through 14 of the preceding section asked us to provide examples of factor rings \BookInline{R/N} where \BookInline{R} and \BookInline{R/N} have very different structural properties. We start with some examples of this situation, and in the process, provide solutions to those exercises.

\BookLabel{31.1 Example} As was shown in Corollary 23.5, the ring \BookInline{\mathbb{Z}_{p},\allowbreak } which is isomorphic to \BookInline{\mathbb{Z}/p\mathbb{Z},\allowbreak } is a field for \BookInline{p} a prime.\textit{ Thus} \BookInline{a} \textit{factor ring of an integral domain may be} \BookInline{a\text{fi}eld.}\BookEnd{\blacktriangle}

\BookLabel{31.2 Example} The ring \BookInline{\mathbb{Z}\times \mathbb{Z}} is not an integral domain, for

\begin{center}\begin{adjustbox}{max width=\linewidth}\BookInline{\displaystyle (0,\allowbreak 1)(1,\allowbreak 0)= (0,\allowbreak 0),\allowbreak }\end{adjustbox}\end{center}

\noindent showing that (0, 1) and (1, 0) are 0 divisors. Let \BookInline{N= \{(0,\allowbreak n)\mid n\in \mathbb{Z}\}.} Now \BookInline{N} is an ideal of \BookInline{\mathbb{Z}\times \mathbb{Z},\allowbreak } and \BookInline{(\mathbb{Z}\times \mathbb{Z})/N} is isomorphic to \BookInline{\mathbb{Z}} under the correspondence \BookInline{[(m,\allowbreak 0)+ N]\leftrightarrow } \BookInline{m,\allowbreak } where \BookInline{m\in \mathbb{Z}.} Thus a\textit{ factor} ring of a ring may be an integral domain, even though the original ring is not\BookInline{.}\BookEnd{\blacktriangle}

\BookLabel{31.3 Example} The subset \BookInline{N= \{0,\allowbreak 3\}} of \BookInline{\mathbb{Z}_{6}} is easily seen to be an ideal of \BookInline{\mathbb{Z}_{6},\allowbreak } and \BookInline{\mathbb{Z}_{6}/N} has three elements\BookInline{,\allowbreak 0+ N,\allowbreak 1+ N,\allowbreak } and \BookInline{2+ N.} These add and multiply in such a fashion as to show that \BookInline{\mathbb{Z}_{6}/N\simeq \mathbb{Z}_{3}} under the correspondence

\begin{center}\begin{adjustbox}{max width=\linewidth}\BookInline{\displaystyle (0+ N)\leftrightarrow 0,\allowbreak (1+ N)\leftrightarrow 1,\allowbreak (2+ N)\leftrightarrow 2.}\end{adjustbox}\end{center}

\noindent This example shows that\textit{ if} \BookInline{R} \textit{is not even an integral domain, that is, if} \BookInline{R} \textit{has zero} \textit{divisors, it is still possible for} \BookInline{R/N} \textit{to be} \BookInline{a\text{fi}eld.}\BookEnd{\blacktriangle}

\BookLabel{31.4 Example} Note that \BookInline{\mathbb{Z}} is an integral domain, but \BookInline{\mathbb{Z}/6\mathbb{Z}\simeq \mathbb{Z}_{6}} is not. The preceding examples showed that a factor ring may have a structure that seems\textit{ better} than the original ring.

\end{BookBody}

% Source PDF page 268; original printed page 251
\BookSource{268}{251}
\BookNarrativetrue

\begin{BookBody}

\noindent This example indicates that the structure of a factor ring may seem\textit{ worse} than that of the original ring\BookInline{.}\BookEnd{\blacktriangle}

Every nonzero ring \BookInline{R} has at least two ideals, the\textbf{ improper ideal} \BookInline{R} and the\textbf{ trivial} \textbf{ideal} \BookInline{\{0\}.} For these ideals, the factor rings are \BookInline{R/R,\allowbreak } which has only one element, and \BookInline{R/\{0\},\allowbreak } which is isomorphic to \BookInline{R.} These are uninteresting cases. Just as for a subgroup of a group, a\textbf{ proper nontrivial ideal} of a ring \BookInline{R} is an ideal \BookInline{N} of \BookInline{R} such that \BookInline{N\neq R} and \BookInline{N\neq \{0\}.}

While factor rings of rings and integral domains may be of great interest, as the above examples indicate, Corollary 31.6, which follows our next theorem, shows that a factor ring of a field is really not useful to us.

\BookLabel{31.5 Theorem} If \BookInline{R} is a ring with unity, and \BookInline{N} is an ideal of \BookInline{R} containing a unit, then \BookInline{N= R.}

\BookLabel{\textit{Proof}} Let \BookInline{N} be an ideal of \BookInline{R,\allowbreak } and suppose that \BookInline{u\in N} for some unit \BookInline{u} in \BookInline{R.} Then the condition \BookInline{rN\subseteq N} for all \BookInline{r\in R} implies, if we take \BookInline{r= u^{- 1}} and \BookInline{u\in N,\allowbreak } that \BookInline{1= u^{- 1}u} is in \BookInline{N.} But then \BookInline{rN\subseteq N} for all \BookInline{r\in R} implies that \BookInline{r1= r} is in \BookInline{N} for all \BookInline{r\in R,\allowbreak } so \BookInline{N= R.}\BookEnd{\blacklozenge}

\BookLabel{31.6 Corollary} A field contains no proper nontrivial ideals.

\BookLabel{\textit{Proof}} Since every nonzero element of a field is a unit, it follows at once from Theorem 31.5 that an ideal of a field \BookInline{F} is either \BookInline{\{0\}} or all of \BookInline{F.}\BookEnd{\blacklozenge}

\BookTextHeading{\textbf{Maximal and Prime Ideals}  }

\noindent We now consider the questions of when a factor ring is a field and when it is an integral domain. In our analogy between groups and rings, we noticed that ideals in rings correspond to normal subgroups. Corollary 31.6 states that a field contains no proper nontrivial ideals. In group theory, this corresponds to a group having no proper nontrivial normal subgroups, that is, a simple group. Theorem 13.20 states that a factor group \BookInline{G/H} is simple if and only if \BookInline{H} is a maximal normal subgroup of \BookInline{G.} The following definition is analogous to maximal normal subgroups.

\BookLabel{31.7 Definition} A\textbf{ maximal ideal of a ring} \BookInline{R} is an ideal \BookInline{M} different from \BookInline{R} such that there is no proper ideal \BookInline{N} of \BookInline{R} properly containing \BookInline{M.}\BookEnd{\blacksquare}

\BookLabel{31.8 Example} Let \BookInline{p} be a prime positive integer. We know that \BookInline{\mathbb{Z}/p\mathbb{Z}} is isomorphic to \BookInline{\mathbb{Z}_{p}.} Forgetting about multiplication for the moment and regarding \BookInline{\mathbb{Z}/p\mathbb{Z}} and \BookInline{\mathbb{Z}_{p}} as additive groups, we know that \BookInline{\mathbb{Z}_{p}} is a simple group, and consequently \BookInline{p\mathbb{Z}} must be a maximal normal subgroup of \BookInline{\mathbb{Z}} by Theorem 13.20. Since \BookInline{\mathbb{Z}} is an abelian group and every subgroup is a normal subgroup, we see that \BookInline{p\mathbb{Z}} is a maximal proper subgroup of \BookInline{\mathbb{Z}.} Since \BookInline{p\mathbb{Z}} is an ideal of the ring \BookInline{\mathbb{Z},\allowbreak } it follows that \BookInline{p\mathbb{Z}} is a maximal ideal of \BookInline{\mathbb{Z}.} We know that \BookInline{\mathbb{Z}/p\mathbb{Z}} is isomorphic to the ring \BookInline{\mathbb{Z}_{p},\allowbreak } and that \BookInline{\mathbb{Z}_{p}} is actually a field. Thus \BookInline{\mathbb{Z}/p\mathbb{Z}} is a field. This illustrates the next theorem\BookInline{.}\BookEnd{\blacktriangle}

\BookLabel{31.9 Theorem} \textbf{(Analogue of Theorem 13.20)} Let \BookInline{R} be a commutative ring with unity. Then \BookInline{M} is a maximal ideal of \BookInline{R} if and only if \BookInline{R/M} is a field.

\BookLabel{\textit{Proof}} We first assume that \BookInline{M} is a maximal ideal in \BookInline{R.} Since \BookInline{R} is a commutative ring with unity, so is \BookInline{R/M.} Furthermore, since \BookInline{M\neq R,\allowbreak 0+ M\neq 1+ M} and \BookInline{R/M} is a nonzero ring. Let \BookInline{(a+ M)\in R/M,\allowbreak } with \BookInline{a\notin M,\allowbreak } so that \BookInline{a+ M} is not the additive identity element of \BookInline{R/M.} Suppose \BookInline{a+ M} has no multiplicative inverse in \BookInline{R/M.} Then the set \BookInline{(R/M)(a+ M)= \{(r+ M)(a+ M)\mid (r+ M)\in R/M\}} does not contain \BookInline{1+ M.} We easily see that \BookInline{(R/M)(a+ M)} is an ideal of \BookInline{R/M.} It is nontrivial because \BookInline{a\notin M,\allowbreak } and it is a 
% Source PDF page 269; original printed page 252
\BookSource{269}{252}
proper ideal because it does not contain \BookInline{1+ M.} By Theorem 30.11, if \BookInline{\gamma:R\to R/M} is the canonical homomorphism, then \BookInline{\gamma ^{- 1}[(R/M)(a+ M)]} is a proper ideal of \BookInline{R} properly containing \BookInline{M.} But this contradicts our assumption that \BookInline{M} is a maximal ideal, so \BookInline{a+ M} must have a multiplicative inverse in \BookInline{R/M.}

\end{BookBody}

\BookNarrativetrue

\begin{BookBody}

Conversely, suppose that \BookInline{R/M} is a field. By Theorem 30.11, if \BookInline{N} is any ideal of \BookInline{R} such that \BookInline{M\subset N\subset R} and \BookInline{\gamma } is the canonical homomorphism of \BookInline{R} onto \BookInline{R/M,\allowbreak } then \BookInline{\gamma [N]} is an ideal of \BookInline{R/M} with \BookInline{\{(0+ M)\}\subset \gamma [N]\subset R/M.} But this is contrary to Corollary 31.6, which states that the field \BookInline{R/M} contains no proper nontrivial ideals. Hence if \BookInline{R/M} is a field, then \BookInline{M} is maximal\BookInline{.}\BookEnd{\blacklozenge}

\BookLabel{31.10 Example} Since \BookInline{\mathbb{Z}/n\mathbb{Z}} is isomorphic to \BookInline{\mathbb{Z}_{n}} and \BookInline{\mathbb{Z}_{n}} is a field if and only if \BookInline{n} is a prime, we see that the maximal ideals of \BookInline{\mathbb{Z}} are precisely the ideals \BookInline{p\mathbb{Z}} for prime positive integers \BookInline{p.}\BookEnd{\blacktriangle}

\BookLabel{31.11 Corollary} A commutative ring with unity is a field if and only if it has no proper nontrivial ideals.

\BookLabel{\textit{Proof}} Corollary 31.6 shows that a field has no proper nontrivial ideals.

Conversely, if a commutative ring \BookInline{R} with unity has no proper nontrivial ideals, then \BookInline{\{0\}} is a maximal ideal and \BookInline{R/\{0\},\allowbreak } which is isomorphic to \BookInline{R,\allowbreak } is a field by Theorem 31.9. \BookEnd{\blacklozenge}

We now turn to the question of characterizing, for a commutative ring \BookInline{R} with unity, the ideals \BookInline{N\neq R} such that \BookInline{R/N} is an integral domain. The answer here is rather obvious. The factor ring \BookInline{R/N} will be an integral domain if and only if \BookInline{(a+ N)(b+ N)= N} implies that either

\begin{center}\begin{adjustbox}{max width=\linewidth}\BookInline{\displaystyle a+ N= N\;\text{or}\;b+ N= N.}\end{adjustbox}\end{center}

\noindent This is exactly the statement that \BookInline{R/N} has no divisors of 0, since the coset \BookInline{N} plays the role of 0 in \BookInline{R/N.} Looking at representatives, we see that this condition amounts to saying that \BookInline{ab\in N} implies that either \BookInline{a\in N} or \BookInline{b\in N.}

\BookLabel{31.12 Example} All ideals of \BookInline{\mathbb{Z}} are of the form \BookInline{n\mathbb{Z}.} For \BookInline{n= 0,\allowbreak } we have \BookInline{n\mathbb{Z}= \{0\},\allowbreak } and \BookInline{\mathbb{Z}/\{0\}\simeq \mathbb{Z},\allowbreak } which is an integral domain. For \BookInline{n>0,\allowbreak } we have \BookInline{\mathbb{Z}/n\mathbb{Z}\simeq \mathbb{Z}_{n}} and \BookInline{\mathbb{Z}_{n}} is an integral domain if and only if \BookInline{n} is a prime. Thus the nonzero ideals \BookInline{n\mathbb{Z}} such that \BookInline{\mathbb{Z}/n\mathbb{Z}} is an integral domain are of the form \BookInline{p\mathbb{Z},\allowbreak } where \BookInline{p} is a prime. Of course\BookInline{,\allowbreak \mathbb{Z}/p\mathbb{Z}} is actually a field, so that \BookInline{p\mathbb{Z}} is a maximal ideal of \BookInline{\mathbb{Z}.} Note that for a product \BookInline{rs} of integers to be in \BookInline{p\mathbb{Z},\allowbreak } the prime \BookInline{p} must divide either \BookInline{r} or \BookInline{s.} The role of prime integers in this example makes the use of the word\textit{ prime} in the next definition more reasonable\BookInline{.}\BookEnd{\blacktriangle}

\BookLabel{31.13 Definition} An ideal \BookInline{N\neq R} in a commutative ring \BookInline{R} is a\textbf{ prime ideal} if \BookInline{ab\in N} implies that either \BookInline{a\in N} or \BookInline{b\in N} for \BookInline{a,\allowbreak b\in R.}\BookEnd{\blacksquare}

Note that \BookInline{\{0\}} is a prime ideal in \BookInline{\mathbb{Z},\allowbreak } and indeed, in any integral domain.

\BookLabel{31.14 Example} Note that \BookInline{\mathbb{Z}\times \{0\}} is a prime ideal of \BookInline{\mathbb{Z}\times \mathbb{Z},\allowbreak } for if \BookInline{(a,\allowbreak b)(c,\allowbreak d)\in \mathbb{Z}\times \{0\},\allowbreak } then we must have \BookInline{bd= 0} in \BookInline{\mathbb{Z}.} This implies that either \BookInline{b= 0} so \BookInline{(a,\allowbreak b)\in \mathbb{Z}\times \{0\}} or \BookInline{d= 0} so \BookInline{(c,\allowbreak d)\in } \BookInline{\mathbb{Z}\times \{0\}.} Note that \BookInline{(\mathbb{Z}\times \mathbb{Z})/(\mathbb{Z}\times \{0\})} is isomorphic to \BookInline{\mathbb{Z},\allowbreak } which is an integral domain. \BookEnd{\blacktriangle}

Our remarks preceding Example 31.12 constitute a proof of the following theorem, which is illustrated by Example 31.14.

\BookLabel{31.15 Theorem} Let \BookInline{R} be a commutative ring with unity, and let \BookInline{N\neq R} be an ideal in \BookInline{R.} Then \BookInline{R/N} is an integral domain if and only if \BookInline{N} is a prime ideal in \BookInline{R.}

\end{BookBody}

% Source PDF page 270; original printed page 253
\BookSource{270}{253}
\BookNarrativetrue

\begin{BookBody}

\BookLabel{31.16 Corollary} Every maximal ideal in a commutative ring \BookInline{R} with unity is a prime ideal.

\BookLabel{\textit{Proof}} If \BookInline{M} is maximal in \BookInline{R,\allowbreak } then \BookInline{R/M} is a field, hence an integral domain, and therefore \BookInline{M} is a prime ideal by Theorem \BookInline{31.15.}\BookEnd{\blacklozenge}

The material that has just been presented regarding maximal and prime ideals is very important and we shall be using it quite a lot. We should keep the main ideas well in mind. We must know and understand the definitions of maximal and prime ideals and must remember the following facts that we have demonstrated.

For a commutative ring \BookInline{R} with unity:

\BookExercise{1.}{An ideal \BookInline{M} of \BookInline{R} is maximal if and only if \BookInline{R/M} is a field.}

\BookExercise{2.}{An ideal \BookInline{N} of \BookInline{R} is prime if and only if \BookInline{R/N} is an integral domain.}

\BookExercise{3.}{Every maximal ideal of \BookInline{R} is a prime ideal.}

\BookTextHeading{\textbf{Prime Fields}  }

\noindent We now proceed to show that the rings \BookInline{\mathbb{Z}} and \BookInline{\mathbb{Z}_{n}} form foundations upon which all rings with unity rest, and that \BookInline{\mathbb{Q}} and \BookInline{\mathbb{Z}_{p}} perform a similar service for all fields. Let \BookInline{R} be any ring with unity 1. Recall that by \BookInline{n\cdot 1} we mean \BookInline{1+ 1+ \cdots + 1} for \BookInline{n} summands for \BookInline{n>0,\allowbreak } and \BookInline{(- 1)+ (- 1)+ \cdots + (- 1)} for \BookInline{|n|} summands for \BookInline{n<0,\allowbreak } while \BookInline{n\cdot 1= 0} for \BookInline{n= 0.}

\BookLabel{31.17 Theorem} If \BookInline{R} is a ring with unity 1, then the map \BookInline{\phi:\mathbb{Z}\to R} given by

\begin{center}\begin{adjustbox}{max width=\linewidth}\BookInline{\displaystyle \phi (n)= n\cdot 1}\end{adjustbox}\end{center}

\noindent for \BookInline{n\in \mathbb{Z}} is a homomorphism of \BookInline{\mathbb{Z}} into \BookInline{R.}

\BookLabel{\textit{Proof}} Observe that

\begin{center}\begin{adjustbox}{max width=\linewidth}\BookInline{\displaystyle \phi (n+ m)= (n+ m)\cdot 1= (n\cdot 1)+ (m\cdot 1)= \phi (n)+ \phi (m).}\end{adjustbox}\end{center}

\noindent The distributive laws in \BookInline{R} show that

\begin{center}\begin{adjustbox}{max width=\linewidth}\BookInline{\displaystyle \underbrace { ( 1 + 1 + \cdots + 1 ) } _ { n \mathrm { \scriptsize ~ s u m m a n d s } } \underbrace { ( 1 + 1 + \cdots + 1 ) } _ { m \mathrm { \scriptsize ~ s u m m a n d s } } = \underbrace { ( 1 + 1 + \cdots + 1 ) } _ { n m \mathrm { \scriptsize ~ s u m m a n d s } } .}\end{adjustbox}\end{center}

Thus \BookInline{(n\cdot 1)(m\cdot 1)= (nm)\cdot 1} for \BookInline{n,\allowbreak m>0.} Similar arguments with the distributive laws show that for all \BookInline{n,\allowbreak m\in \mathbb{Z},\allowbreak } we have

\begin{center}\begin{adjustbox}{max width=\linewidth}\BookInline{\displaystyle (n\cdot 1)(m\cdot 1)= (nm)\cdot 1.}\end{adjustbox}\end{center}

\noindent Thus

\begin{center}\begin{adjustbox}{max width=\linewidth}\BookInline{\displaystyle \phi (nm)= (nm)\cdot 1= (n\cdot 1)(m\cdot 1)= \phi (n)\phi (m).}\end{adjustbox}\end{center} \BookEnd{\blacklozenge}

\BookLabel{31.18 Corollary} Let \BookInline{R} be a ring with unity. If \BookInline{R} has characteristic \BookInline{n>1,\allowbreak } then \BookInline{R} contains a subring isomorphic to \BookInline{\mathbb{Z}_{n}.} If \BookInline{R} has characteristic 0, then \BookInline{R} contains a subring isomorphic to \BookInline{\mathbb{Z}.}

\BookLabel{\textit{Proof}} The map \BookInline{\phi:\mathbb{Z}\to R} given by \BookInline{\phi (m)= m\cdot 1} for \BookInline{m\in \mathbb{Z}} is a homomorphism by Theorem 31.17. The kernel must be an ideal in \BookInline{\mathbb{Z}.} All ideals in \BookInline{\mathbb{Z}} are of the form \BookInline{s\mathbb{Z}} for some \BookInline{s\in \mathbb{Z}.} By Theorem 23.14 we see that if \BookInline{R} has characteristic \BookInline{n>0,\allowbreak } then the kernel of \BookInline{\phi } is \BookInline{n\mathbb{Z}.} Then the image \BookInline{\phi [\mathbb{Z}]\leq R} is isomorphic to \BookInline{\mathbb{Z}/n\mathbb{Z}\simeq \mathbb{Z}_{n}.} If the characteristic of \BookInline{R} is 0, then \BookInline{m\cdot 1\neq 0} for all \BookInline{m\neq 0,\allowbreak } so the kernel of \BookInline{\phi } is \BookInline{\{0\}.} Thus, the image \BookInline{\phi [\mathbb{Z}]\leq R} is isomorphic to \BookInline{\mathbb{Z}.}\BookEnd{\blacklozenge}

\end{BookBody}

% Source PDF page 271; original printed page 254
\BookSource{271}{254}
\BookNarrativetrue

\begin{BookBody}

\BookLabel{31.19 Theorem} A field \BookInline{F} is either of prime characteristic \BookInline{p} and contains a subfield isomorphic to \BookInline{\mathbb{Z}_{p}} or of characteristic 0 and contains a subfield isomorphic to \BookInline{\mathbb{Q}.}

\BookLabel{\textit{Proof}} If the characteristic of \BookInline{F} is not 0, the above corollary shows that \BookInline{F} contains a subring isomorphic to \BookInline{\mathbb{Z}_{n}.} Then \BookInline{n} must be a prime \BookInline{p,\allowbreak } or \BookInline{F} would have 0 divisors. If \BookInline{F} is of characteristic 0, then \BookInline{F} must contain a subring isomorphic to \BookInline{\mathbb{Z}.} In this case Corollaries 26.9 and 26.10 show that \BookInline{F} must contain a field of quotients of this subring and that this field of quotients must be isomorphic to \BookInline{\mathbb{Q}.}\BookEnd{\blacklozenge}

Thus every field contains either a subfield isomorphic to \BookInline{\mathbb{Z}_{p}} for some prime \BookInline{p} or a subfield isomorphic to \BookInline{\mathbb{Q}.} These fields \BookInline{\mathbb{Z}_{p}} and \BookInline{\mathbb{Q}} are the fundamental building blocks on which all fields rest.

\BookLabel{31.20 Definition} The fields \BookInline{\mathbb{Z}_{p}} and \BookInline{\mathbb{Q}} are\textbf{ prime fields.} \BookEnd{\blacksquare}

\BookTextHeading{\textbf{Ideal Structure in} \textit{\textbf{F}}\textbf{[}\textit{\textbf{x}}\textbf{]}   }

\noindent Throughout the rest of this section, we assume that \BookInline{F} is a field. We give the next definition for a general commutative ring \BookInline{R} with unity, although we are only interested in the case \BookInline{R= F[x].} Note that for a commutative ring \BookInline{R} with unity and \BookInline{a\in R,\allowbreak } the set \BookInline{\{}\textit{ra} \BookInline{\mid r\in R\}} is an ideal in \BookInline{R} that contains the element \BookInline{a.}

\BookLabel{31.21 Definition} If \BookInline{R} is a commutative ring with unity and \BookInline{a\in R,\allowbreak } the ideal \BookInline{\{}\textit{ra} \BookInline{\mid r\in R\}} of all multiples of \BookInline{a} is the\textbf{ principal ideal generated by} \BookInline{a} and is denoted by \BookInline{\langle a\rangle.} An ideal \BookInline{N} of \BookInline{R} is a \textbf{principal ideal} if \BookInline{N= \langle a\rangle } for some \BookInline{a\in R.}\BookEnd{\blacksquare}

\BookLabel{31.22 Example} Every ideal of the ring \BookInline{\mathbb{Z}} is of the form \BookInline{n\mathbb{Z},\allowbreak } which is generated by \BookInline{n,\allowbreak } so every ideal of \BookInline{\mathbb{Z}} is a principal ideal\BookInline{.}\BookEnd{\blacktriangle}

\BookLabel{31.23 Example} The principal ideal \BookInline{\langle x\rangle } in \BookInline{F[x]} consists of all polynomials in \BookInline{F[x]} having zero constant term\BookInline{.}\BookEnd{\blacktriangle}

The next theorem is another simple but very important application of the division algorithm for \BookInline{F[x].(}See Theorem 28.2.) The proof of this theorem is to the division algorithm in \BookInline{F[x]} as the proof that a subgroup of a cyclic group is cyclic is to the division algorithm in \BookInline{\mathbb{Z}.}

\BookLabel{31.24 Theorem} If \BookInline{F} is a field, every ideal in \BookInline{F[x]} is principal.

\BookLabel{\textit{Proof}} Let \BookInline{N} be an ideal of \BookInline{F[x].} If \BookInline{N= \{0\},\allowbreak } then \BookInline{N= \langle 0\rangle.} Suppose that \BookInline{N\neq \{0\},\allowbreak } and let \BookInline{g(x)} be a nonzero element of \BookInline{N} of minimal degree. If the degree of \BookInline{g(x)} is 0, then \BookInline{g(x)\in F} and is a unit, so \BookInline{N= F[x]= \langle 1\rangle } by Theorem 31.5, so \BookInline{N} is principal. If the degree of \BookInline{g(x)} is \BookInline{\geq 1,\allowbreak } let \BookInline{f(x)} be any element of \BookInline{N.} Then by Theorem \BookInline{28.2,\allowbreak f(x)= g(x)q(x)+ r(x),\allowbreak } where \BookInline{r(x)= 0} or (degree \BookInline{r(x))<(}degree \BookInline{g(x)).} Now \BookInline{f(x)\in N} and \BookInline{g(x)\in N} imply that \BookInline{f(x)- g(x)q(x)= r(x)} is in \BookInline{N} by definition of an ideal. Since \BookInline{g(x)} is a nonzero element of minimal degree in \BookInline{N,\allowbreak } we must have \BookInline{r(x)= 0.} Thus \BookInline{f(x)= g(x)q(x)} and \BookInline{N= \langle g(x)\rangle.} \BookEnd{\blacklozenge}

We can now characterize the maximal ideals of \BookInline{F[x].} This is a crucial step in achieving our\textbf{ basic goal:} to show that any nonconstant polynomial \BookInline{f(x)} in \BookInline{F[x]} has a zero in some field \BookInline{E} containing \BookInline{F.}

\BookLabel{31.25 Theorem} An ideal \BookInline{\langle p(x)\rangle \neq \{0\}} of \BookInline{F[x]} is maximal if and only if \BookInline{p(x)} is irreducible over \BookInline{F.}

\end{BookBody}

% Source PDF page 272; original printed page 255
\BookSource{272}{255}
\BookNarrativetrue

\begin{BookBody}

\BookLabel{\textit{Proof}} Suppose that \BookInline{\langle p(x)\rangle \neq \{0\}} is a maximal ideal of \BookInline{F[x].} Then \BookInline{\langle p(x)\rangle \neq F[x],\allowbreak } so \BookInline{p(x)\not \in F.} Let \BookInline{p(x)= f(x)g(x)} be a factorization of \BookInline{p(x)} in \BookInline{F[x].} Since \BookInline{\langle p(x)\rangle } is a maximal ideal and hence also a prime ideal\BookInline{,\allowbreak (f(x)g(x))\in \langle p(x)\rangle } implies that \BookInline{f(x)\in \langle p(x)\rangle } or \BookInline{g(x)\in \langle p(x)\rangle;} that is, either \BookInline{f(x)} or \BookInline{g(x)} has \BookInline{p(x)} as a factor. But then we can’t have the degrees of both \BookInline{f(x)} and \BookInline{g(x)} less than the degree of \BookInline{p(x).} This shows that \BookInline{p(x)} is irreducible over \BookInline{F.}

Conversely, if \BookInline{p(x)} is irreducible over \BookInline{F,\allowbreak } suppose that \BookInline{N} is an ideal such that \BookInline{\langle p(x)\rangle \subseteq N\subseteq F[x].} Now \BookInline{N} is a principal ideal by Theorem 31.24, so \BookInline{N= \langle g(x)\rangle } for some \BookInline{g(x)\in N.} Then \BookInline{p(x)\in N} implies that \BookInline{p(x)= g(x)q(x)} for some \BookInline{q(x)\in F[x].} But \BookInline{p(x)} is irreducible, which implies that either \BookInline{g(x)} or \BookInline{q(x)} is of degree 0. If \BookInline{g(x)} is of degree 0, that is, a nonzero constant in \BookInline{F,\allowbreak } then \BookInline{g(x)} is a unit in \BookInline{F[x],\allowbreak } so \BookInline{\langle g(x)\rangle = N= F[x].} If \BookInline{q(x)} is of degree 0, then \BookInline{q(x)= c,\allowbreak } where \BookInline{c\in F,\allowbreak } and \BookInline{g(x)= (1/c)p(x)} is in \BookInline{\langle p(x)\rangle,\allowbreak } so \BookInline{N= \langle p(x)\rangle.} Thus \BookInline{\langle p(x)\rangle \subset N\subset F[x]} is impossible, so \BookInline{\langle p(x)\rangle } is maximal\BookInline{.}\BookEnd{\blacklozenge}

\BookLabel{31.26 Example} Example 28.10 shows that \BookInline{x^{3}+ 3x+ 2} is irreducible in \BookInline{\mathbb{Z}_{5}[x],\allowbreak } so \BookInline{\mathbb{Z}_{5}[x]/\langle x^{3}+ 3x+ 2\rangle } is a field. Similarly, Theorem 27.11 shows that \BookInline{x^{2}- 2} is irreducible in \BookInline{\mathbb{Q}[x],\allowbreak } so \BookInline{\mathbb{Q}[x]/} \BookInline{\langle x^{2}- 2\rangle } is a field. We shall examine such fields in more detail later\BookInline{.}\BookEnd{\blacktriangle}

\BookTextHeading{\textbf{Application to Unique Factorization in} \textit{\textbf{F}}\textbf{[}\textit{\textbf{x}}\textbf{]}   }

\noindent In Section 28, we stated without proof Theorem 31.27, which follows. (See Theorem 28.19.) Assuming this theorem, we proved in Section 28 that factorization of polynomials in \BookInline{F[x]} into irreducible polynomials is unique, except for order of factors and units in \BookInline{F.} We delayed the proof of Theorem 31.27 until now since the machinery we have developed enables us to give such a simple proof. This proof fills the gap in our proof of unique factorization in \BookInline{F[x].}

\BookLabel{31.27 Theorem} Let \BookInline{p(x)} be an irreducible polynomial in \BookInline{F[x].} If \BookInline{p(x)} divides \BookInline{r(x)s(x)} for \BookInline{r(x),\allowbreak s(x)\in } \BookInline{F[x],\allowbreak } then either \BookInline{p(x)} divides \BookInline{r(x)} or \BookInline{p(x)} divides \BookInline{s(x).}

\BookLabel{\textit{Proof}} Suppose \BookInline{p(x)} divides \BookInline{r(x)s(x).} Then \BookInline{r(x)s(x)\in \langle p(x)\rangle,\allowbreak } which is maximal by Theorem 31.25. Therefore\BookInline{,\allowbreak \langle p(x)\rangle } is a prime ideal by Corollary 31.16. Hence \BookInline{r(x)s(x)\in \langle p(x)\rangle } implies that either \BookInline{r(x)\in \langle p(x)\rangle,\allowbreak } giving \BookInline{p(x)} divides \BookInline{r(x),\allowbreak } or that \BookInline{s(x)\in \langle p(x)\rangle,\allowbreak } giving \BookInline{p(x)} divides \BookInline{s(x).}\BookEnd{\blacklozenge}

\BookTextHeading{\textbf{A Preview of Our Basic Goal}  }

\noindent We close this section with an outline of the demonstration in Section 39 of our basic goal. We have all the ideas for the proof at hand now; perhaps you can fill in the details from this outline.

\textbf{Basic goal:} Let \BookInline{F} be a field and let \BookInline{f(x)} be a nonconstant polynomial in \BookInline{F[x].} Show that there exists a field \BookInline{E} containing \BookInline{F} and containing a zero \BookInline{\alpha } of \BookInline{f(x).}

\BookTextHeading{\textbf{Outline of the Proof}  }

\BookExercise{1.}{Let \BookInline{p(x)} be an irreducible factor of \BookInline{f(x)} in \BookInline{F[x].}}

\BookExercise{2.}{Let \BookInline{E} be the \BookInline{\text{fi}eldF[x]/\langle p(x)\rangle.(}See Theorems 31.25 and 31.9.)}

\BookExercise{3.}{Show that no two different elements of \BookInline{F} are in the same coset of \BookInline{F[x]/\langle p(x)\rangle,\allowbreak } and deduce that we may consider \BookInline{F} to be (isomorphic to) a subfield of \BookInline{E.}}

\BookExercise{4.}{Let \BookInline{\alpha } be the coset \BookInline{x+ \langle p(x)\rangle } in \BookInline{E.} Show that for the evaluation homomorphism \BookInline{\phi _{\alpha }:F[x]\to E,\allowbreak } we have \BookInline{\phi _{\alpha }(f(x))= 0.} That is\BookInline{,\allowbreak \alpha } is a zero of \BookInline{f(x)} in \BookInline{E.}}

\end{BookBody}

% Source PDF page 273; original printed page 256
\BookSource{273}{256}
\BookNarrativetrue

\begin{BookBody}

An example of a field constructed according to this outline is given in Section 39. There, we give addition and multiplication tables for the field \BookInline{\mathbb{Z}_{2}[x]/\langle x^{2}+ x+ 1\rangle.} We show there that this field has just four elements, the cosets

\begin{center}\begin{adjustbox}{max width=\linewidth}\BookInline{\displaystyle 0+ \langle x^{2}+ x+ 1\rangle ,\allowbreak 1+ \langle x^{2}+ x+ 1\rangle ,\allowbreak x+ \langle x^{2}+ x+ 1\rangle ,\allowbreak }\end{adjustbox}\end{center}

\noindent and

\begin{center}\begin{adjustbox}{max width=\linewidth}\BookInline{\displaystyle (x+ 1)+ \langle x^{2}+ x+ 1\rangle .}\end{adjustbox}\end{center}

\noindent We rename these four cosets \BookInline{0,\allowbreak 1,\allowbreak \alpha,\allowbreak } and \BookInline{\alpha + 1} respectively, and obtain Tables 39.21 and 39.22 for addition and multiplication in this 4-element field. To see how these tables are constructed, remember that we are in a field of characteristic 2, so that \BookInline{\alpha + \alpha = } \BookInline{\alpha (1+ 1)= \alpha 0= 0.} Remember also that \BookInline{\alpha } is a zero of \BookInline{x^{2}+ x+ 1,\allowbreak } so that \BookInline{\alpha ^{2}+ \alpha + 1= } 0 and consequently \BookInline{\alpha ^{2}= - \alpha - 1= \alpha + 1.}

\end{BookBody}

\BookNarrativefalse

\begin{BookBody}

\BookExerciseTitle{31} \BookEnd{\blacksquare}

\BookTextHeading{\textbf{Computations}}

\BookExercise{1.}{Find all prime ideals and all maximal ideals of \BookInline{\mathbb{Z}_{6}.}}

\BookExercise{2.}{Find all prime ideals and all maximal ideals of \BookInline{\mathbb{Z}_{12}.}}

\BookExercise{3.}{Find all prime ideals and all maximal ideals of \BookInline{\mathbb{Z}_{2}\times \mathbb{Z}_{2}.}}

\BookExercise{4.}{Find all prime ideals and all maximal ideals of \BookInline{\mathbb{Z}_{2}\times \mathbb{Z}_{4}.}}

\BookExercise{5.}{Find all \BookInline{c\in \mathbb{Z}_{3}} such that \BookInline{\mathbb{Z}_{3}[x]/\langle x^{2}+ c\rangle } is a field.}

\BookExercise{6.}{Find all \BookInline{c\in \mathbb{Z}_{3}} such that \BookInline{\mathbb{Z}_{3}[x]/\langle x^{3}+ x^{2}+ c\rangle } is a field.}

\BookExercise{7.}{Find all \BookInline{c\in \mathbb{Z}_{3}} such that \BookInline{\mathbb{Z}_{3}[x]/\langle x^{3}+ cx^{2}+ 1\rangle } is a field.}

\BookExercise{8.}{Find all \BookInline{c\in \mathbb{Z}_{5}} such that \BookInline{\mathbb{Z}_{5}[x]/\langle x^{2}+ x+ c\rangle } is a field.}

\BookExercise{9.}{Find all \BookInline{c\in \mathbb{Z}_{5}} such that \BookInline{\mathbb{Z}_{5}[x]/\langle x^{2}+ cx+ 1\rangle } is a field.}

\BookTextHeading{\textbf{Concepts}}

\noindent In Exercises 10 through 13, correct the definition of the italicized term without reference to the text, if correction is needed, so that it is in a form acceptable for publication.

\BookExercise{10.}{A\textit{ maximal ideal} of a ring \BookInline{R} is an ideal that is not contained in any other ideal of \BookInline{R.}}

\BookExercise{11.}{A\textit{ prime ideal} of a commutative ring \BookInline{R} is an ideal of the form \BookInline{pR= \{pr\mid r\in R\}} for some prime \BookInline{p.}}

\BookExercise{12.}{A\textit{ prime} \BookInline{\text{fi}eld} is a field that has no proper subfields.}

\BookExercise{13.}{A\textit{ principal ideal} of a commutative ring with unity is an ideal \BookInline{N} with the property that there exists \BookInline{a\in N} such that \BookInline{N} is the smallest ideal that contains \BookInline{a.}}

\BookExercise{14.}{Determine whether each of the following is true or false.}

\BookExercise{a.}{Every prime ideal of every commutative ring with unity is a maximal ideal.}

\BookExercise{b.}{Every maximal ideal of every commutative ring with unity is a prime ideal.}

\BookExercise{c.}{\textbf{c}\BookInline{.\mathbb{Q}} is its own prime subfield.}

\BookExercise{d.}{The prime subfield of \BookInline{\mathbb{C}} is \BookInline{\mathbb{R}.}}

\BookExercise{e.}{Every field contains a subfield isomorphic to a prime field.}

\BookExercise{f.}{A ring with zero divisors may contain one of the prime fields as a subring.}

\BookExercise{g.}{Every field of characteristic zero contains a subfield isomorphic to \BookInline{\mathbb{Q}.}}

\BookExercise{h.}{Let \BookInline{F} be a field. Since \BookInline{F[x]} has no divisors of 0, every ideal of \BookInline{F[x]} is a prime ideal.}

\BookExercise{i.}{Let \BookInline{F} be a field. Every ideal of \BookInline{F[x]} is a principal ideal.}

\BookExercise{j.}{Let \BookInline{F} be a field. Every principal ideal of \BookInline{F[x]} is a maximal ideal.}

\BookExercise{15.}{Find a maximal ideal of \BookInline{\mathbb{Z}\times \mathbb{Z}.}}

\end{BookBody}

% Source PDF page 274; original printed page 257
\BookSource{274}{257}
\BookNarrativefalse

\begin{BookBody}

\BookExercise{16.}{Find a prime ideal of \BookInline{\mathbb{Z}\times \mathbb{Z}} that is not maximal.}

\BookExercise{17.}{Find a nontrivial proper ideal of \BookInline{\mathbb{Z}\times \mathbb{Z}} that is not prime.}

\BookExercise{18.}{Is \BookInline{\mathbb{Q}[x]/\langle x^{2}- 5x+ 6\rangle } a field? Why?}

\BookExercise{19.}{Is \BookInline{\mathbb{Q}[x]/\langle x^{2}- 6x+ 6\rangle } a field? Why?}

\BookLabel{\textit{Proof}} \textbf{ Synopsis}

\BookExercise{20.}{Give a one- or two-sentence synopsis of “only if” part of Theorem 31.9.}

\BookExercise{21.}{Give a one- or two-sentence synopsis of “if” part of Theorem 31.9.}

\BookExercise{22.}{Give a one- or two-sentence synopsis of Theorem 31.24.}

\BookExercise{23.}{Give a one- or two-sentence synopsis of the “only if” part of Theorem 31.25.}

\BookTextHeading{\textbf{Theory}}

\BookExercise{24.}{Give an example of an ideal in \BookInline{\mathbb{Q}[x,\allowbreak y]} that is not a principal ideal. Conclude that if \BookInline{R} is an integral domain with the property that every ideal in \BookInline{R} is principal, it does not follow that every ideal in \BookInline{R[x]} is a principal ideal.}

\BookExercise{25.}{Prove that if \BookInline{R} is a commutative ring with unity and \BookInline{a\in R,\allowbreak } then \BookInline{\langle a\rangle = \{}\textit{ra} \BookInline{\mid r\in R\}} is an ideal in \BookInline{R.}}

\BookExercise{26.}{Let \BookInline{R} be a finite commutative ring with unity. Show that every prime ideal in \BookInline{R} is a maximal ideal.}

\BookExercise{27.}{Corollary 31.18 tells us that every ring with unity contains a subring isomorphic to either \BookInline{\mathbb{Z}} or some \BookInline{\mathbb{Z}_{n}.} Is it possible that a ring with unity may simultaneously contain two subrings isomorphic to \BookInline{\mathbb{Z}_{n}} and \BookInline{\mathbb{Z}_{m}} for \BookInline{n\neq m}? If it is possible, give an example. If it is impossible, prove it.}

\BookExercise{28.}{Continuing Exercise 27, is it possible that a ring with unity may simultaneously contain two subrings isomorphic to the fields \BookInline{\mathbb{Z}_{p}} and \BookInline{\mathbb{Z}_{q}} for two different primes \BookInline{p} and \BookInline{q}? Give an example or prove it is impossible.}

\BookExercise{29.}{Following the idea of Exercise 28, is it possible for an integral domain to contain two subrings isomorphic to \BookInline{\mathbb{Z}_{p}} and \BookInline{\mathbb{Z}_{q}} for \BookInline{p\neq q} and \BookInline{p} and \BookInline{q} both prime? Give reasons or an illustration.}

\BookExercise{30.}{Prove directly from the definitions of maximal and prime ideals that every maximal ideal of a commutative ring \BookInline{R} with unity is a prime ideal. [\textit{Hint}: Suppose \BookInline{M} is maximal in \BookInline{R,\allowbreak ab\in M,\allowbreak } and \BookInline{a\not \in M.} Argue that the smallest ideal \BookInline{\{}\textit{ra} \BookInline{+ m\mid r\in R,\allowbreak m\in M\}} containing \BookInline{a} and \BookInline{M} must contain 1. Express 1 as\textit{ ra} \BookInline{+ m} and multiply by \BookInline{b.]}}

\BookExercise{31.}{Show that \BookInline{N} is a maximal ideal in a ring \BookInline{R} if and only if \BookInline{R/N} is a\textbf{ simple ring,} that is, it is nontrivial and has no proper nontrivial ideals. (Compare with Theorem 13.20.)}

\BookExercise{32.}{Prove that if \BookInline{F} is a field, every proper nontrivial prime ideal of \BookInline{F[x]} is maximal.}

\BookExercise{33.}{Let \BookInline{F} be a field and \BookInline{f(x),\allowbreak g(x)\in F[x].} Show that \BookInline{f(x)} divides \BookInline{g(x)} if and only if \BookInline{g(x)\in \langle f(x)\rangle.}}

\BookExercise{34.}{Let \BookInline{F} be a field and let \BookInline{f(x),\allowbreak g(x)\in F[x].} Show that}

\begin{center}\begin{adjustbox}{max width=\linewidth}\BookInline{\displaystyle N= \{r(x)f(x)+ s(x)g(x)\mid r(x),\allowbreak s(x)\in F[x]\}}\end{adjustbox}\end{center}

is an ideal of \BookInline{F[x].} Show that if \BookInline{f(x)} and \BookInline{g(x)} have different degrees and \BookInline{N\neq F[x],\allowbreak } then \BookInline{f(x)} and \BookInline{g(x)} cannot both be irreducible over \BookInline{F.}

\BookExercise{35.}{Use Theorem 31.24 to prove the\textit{ equivalence} of these two theorems:}

\textbf{Fundamental Theorem of Algebra:} Every nonconstant polynomial in \BookInline{\mathbb{C}[x]} has a zero in \BookInline{\mathbb{C}.}

\textbf{Nullstellensatz for} \BookInline{\mathbb{C}[x]:} Let \BookInline{f_{1}(x),\allowbreak \cdots,\allowbreak f_{r}(x)\in \mathbb{C}[x]} and suppose that every \BookInline{\alpha \in \mathbb{C}} that is a zero of all \BookInline{r} of these polynomials is also a zero of a polynomial \BookInline{g(x)} in \BookInline{\mathbb{C}[x].} Then some power of \BookInline{g(x)} is in the smallest ideal of \BookInline{\mathbb{C}[x]} that contains the \BookInline{r} polynomials \BookInline{f_{1}(x),\allowbreak \cdots,\allowbreak f_{r}(x).}

\noindent There is a sort of arithmetic of ideals in a ring. The next three exercises define sum, product, and quotient of ideals.

\BookExercise{36.}{If \BookInline{A} and \BookInline{B} are ideals of a ring \BookInline{R,\allowbreak } the\textbf{ sum} \BookInline{A+ B} \textbf{of} \BookInline{A} \textbf{and} \BookInline{B} is defined by}

\begin{center}\begin{adjustbox}{max width=\linewidth}\BookInline{\displaystyle A+ B= \{a+ b\mid a\in A,\allowbreak b\in B\}.}\end{adjustbox}\end{center}

\BookExercise{a.}{Show that \BookInline{A+ B} is an ideal.}

\BookExercise{b.}{Show that \BookInline{A\subseteq A+ B} and \BookInline{B\subseteq A+ B.}}

\end{BookBody}

% Continuation on source PDF page 275, preceding Section 32
\BookNarrativefalse

\begin{BookBody}

\BookExercise{37.}{Let \BookInline{A} and \BookInline{B} be ideals of a ring \BookInline{R.} The\textbf{ product} \BookInline{AB} \textbf{of} \BookInline{A} \textbf{and} \BookInline{B} is defined by}

\begin{center}\begin{adjustbox}{max width=\linewidth}\BookInline{\displaystyle AB= \Big\{\sum _{i= 1}^{n}a_{i}b_{i}\mid a_{i}\in A,\allowbreak b_{i}\in B,\allowbreak n\in \mathbb{Z}^{+ }\Big\}.}\end{adjustbox}\end{center}

\BookExercise{a.}{Show that \BookInline{AB} is an ideal in \BookInline{R.}}

\BookExercise{b.}{Show that \BookInline{AB\subseteq (A\cap B).}}

\BookExercise{38.}{Let \BookInline{A} and \BookInline{B} be ideals of a\textit{ commutative} ring \BookInline{R.} The\textbf{ quotient} \BookInline{A:B} \textbf{of} \BookInline{A} \textbf{by} \BookInline{B} is defined by}

\begin{center}\begin{adjustbox}{max width=\linewidth}\BookInline{\displaystyle A:B= \{r\in R\mid rb\in A\;\text{for all}\;b\in B\}.}\end{adjustbox}\end{center}

Show that \BookInline{A:B} is an ideal of \BookInline{R.}

\BookExercise{39.}{Show that for a field \BookInline{F,\allowbreak } the set \BookInline{S} of all matrices of the form}

\begin{center}\begin{adjustbox}{max width=\linewidth}\BookInline{\displaystyle \begin{pmatrix} a & b \\ 0 & 0 \end{pmatrix}}\end{adjustbox}\end{center}

for \BookInline{a,\allowbreak b\in F} is a\textbf{ right ideal} but not a\textbf{ left ideal} of \BookInline{M_{2}(F).} That is, show that \BookInline{S} is a subring closed under multiplication on the\textit{ right} by any element of \BookInline{M_{2}(F),\allowbreak } but is not closed under\textit{ left} multiplication.

\BookExercise{40.}{Show that the matrix ring \BookInline{M_{2}(\mathbb{Z}_{2})} is a simple ring; that is\BookInline{,\allowbreak M_{2}(\mathbb{Z}_{2})} has no proper nontrivial ideals.}

\end{BookBody}
\ifdefined\BookMaster\let\BookFinish\relax\else\def\BookFinish{\end{document}}\fi\BookFinish
